\chapter[The Checkpoint Boundary]{The Checkpoint Boundary: Recover Training, Export Inference}
\label{ch:checkpoint-boundary}

A file containing model weights is not necessarily a training checkpoint.
Loading identical weights can reproduce a forward pass while producing a
different next optimizer update. This chapter makes that distinction executable
before introducing distributed checkpointing.

\begin{figure}[htbp]\centering
\begin{tikzpicture}
\node[enginestate,text width=4.6cm] (resume) at (0,0) {RESUME ARTIFACT\\weights, optimizer moments\\step, data identity, RNG\\config, tokenizer identity};
\node[engineop,text width=4.6cm] (export) at (6.5,0) {DEPLOYMENT EXPORT\\weights, architecture contract\\tokenizer, numerical policy\\conversion checks};
\node[enginevalue,text width=4.6cm] (train) at (0,-3.2) {Training continuation\\same next batch and update};
\node[enginevalue,text width=4.6cm] (infer) at (6.5,-3.2) {Inference runtime\\request state, KV memory\\scheduling, output stream};
\draw[enginearrow] (resume) -- node[above,enginenote]{transform} (export);
\draw[enginearrow] (resume) -- node[left,enginenote]{restore} (train);
\draw[enginearrow] (export) -- node[right,enginenote]{load} (infer);
\node[enginenote,text width=11cm] at (3.25,-5) {The arrow across the boundary is not reversible in general.\\A deployment export does not recover discarded optimizer or data state.};
\end{tikzpicture}

\caption{Training and inference meet at a versioned artifact contract. The
export path shown is a design boundary; a cross-book converter is not yet implemented.}
\label{fig:checkpoint-cut}
\end{figure}

\section{Define the Cut Before Saving the File}

Our supported cut is a completed CPU FP32 optimizer step with gradients cleared.
There is no dropout, learning-rate scheduler, mixed-precision scaler, data-loader
worker, distributed rank or in-flight collective. Those omissions make a small
exact-resume claim possible. Adding any of them requires extending the state
contract and its test, not merely adding a new command-line flag.

The schema saves model weights, AdamW state, optimizer step, model configuration,
tokenizer identity, token-stream identity, sampled-batch count, private data RNG
state and global CPU Torch RNG state. Alphabet order participates in the
tokenizer hash: two vocabularies with the same size need not give an ID the same
meaning. Dataset identity includes tokens, context width and vocabulary size.
The file also records the PyTorch version; loading a different version is rejected.

\section{Why Weights Alone Are Insufficient}

Let \(g_t\) be the current gradient and \(m_t,v_t\) AdamW's moment estimates.
The next update depends on all three. Restoring only \(\theta_t\) while setting
\(m_t=v_t=0\) changes the update even if the next batch is identical. Restoring
moments while sampling a different batch changes \(g_t\). Matching data and
moments while changing a stochastic operation's random state can also diverge.

The exact-resume test runs four updates, saves, runs four more, then constructs
a new model and optimizer. It restores the saved state and repeats the same
four updates. Losses, every parameter, optimizer state, batch count and the next
global RNG draw must agree exactly on the tested CPU runtime. Construction
consumes random values, so RNG restoration happens last. This proves a bounded
local continuation, not cross-device or cross-version reproducibility.

\section{Atomic Visibility Is Not Durability}

The writer serializes into a unique temporary file in the destination directory,
closes it, then replaces the destination path. Readers on the same filesystem
see the old file or the new file, not a partially overwritten destination.
This implementation does not synchronize file and directory metadata to durable
storage. Power loss, remote object stores, corrupt disks and distributed commit
protocols remain outside its guarantee.

Only load trusted local checkpoints. The reader uses
\texttt{torch.load(..., weights\_only=True)}, but a restricted deserializer is
not a complete security boundary against malformed or resource-exhausting
artifacts. Identity checks reject wrong tokenizer, data, model configuration,
schema and runtime. They do not authenticate an attacker-controlled file.
Checksums, signatures, access control and quota enforcement are future operations work.

\section{A Memory Budget With Units}

For \(P\) FP32 parameters, a simple persistent training-state payload is
\term{DERIVED}:
\[
 M_{\mathrm{persistent}}\approx
 \underbrace{4P}_{\text{parameters}}+
 \underbrace{4P}_{\text{gradients}}+
 \underbrace{8P}_{\text{two Adam moments}}=16P\ \text{bytes}.
\]
This counts gradients while they exist. It excludes step tensors, allocator
rounding, saved activations, optimizer temporaries, token batches and library
workspaces. After gradients are cleared, the corresponding payload may be
released, though an allocator may retain reserved memory. The unfused attention
scores alone require \(4BH_qT^2\) bytes in FP32. Doubling sequence length
quadruples this particular allocation; it does not prove total memory quadruples.

Mixed-precision policies require a new accounting table: storage dtype,
compute dtype, accumulation dtype, master parameters, gradients and moments
must be stated separately. Do not copy a universal bytes-per-parameter number
into a report without matching the actual implementation.

\section{The Road Toward Distributed Training}

\begin{center}
\begin{tabular}{p{2.1cm}p{4.3cm}p{6cm}}\toprule
Method & Partitioned quantity & Correctness gate before performance \\\midrule
Data parallel & Different examples, replicated model & Match one-process global-batch update \\
State sharding & Optimizer, gradients or parameters & Gather/reshard and resume equivalence \\
Tensor parallel & Operator dimensions & Match local operator and required reductions \\
Pipeline parallel & Layers and microbatch schedule & Correct weight version and gradient schedule \\
Context parallel & Sequence positions & Global attention and position equivalence \\
Expert parallel & Experts and routed activations & Routing, capacity and dispatch/combine equivalence \\\bottomrule
\end{tabular}
\end{center}

These are planned capabilities, not six implemented backends. Megatron-LM and
ZeRO are foundational references for partitioning work and persistent state
\cite{shoeybi2019megatron,rajbhandari2020zero}. The next executable distributed
gate should be two-process data parallel equivalence on a tiny batch, followed
by interrupted recovery. Until those tests exist, throughput comparisons across
multiple devices are premature.

\section{The Companion Book Owns Runtime State}

\emph{Inside the LLM Engine: From model weights and KV memory to industrial
inference serving} owns the detailed treatment of request admission, paged KV
memory, continuous batching, prefix-cache identity, speculative scheduling and
service reliability. This book owns gradients, training data, optimizer state,
distributed training and post-training objectives. Both must agree about model
architecture, tensor layout, tokenizer semantics and numerical tolerances.

A future export must record head counts, RoPE pairing convention and base,
normalization epsilon, vocabulary order, weight tying, linear orientation and
dtype. Validate selected intermediate tensors and full logits before comparing
generated text: sampling can hide or amplify small numerical differences.
GGUF may be a deployment artifact, but it is not the native resume format here.
Exporting GQA weights is also not proof that an inference engine implements
the same head grouping.

\labtitle{Recover the Next Update}
Run \texttt{test\_checkpoint\_exact\_resume\_and\_identity\_rejection}.
Explain every saved field by removing it in a local experiment. Change the
alphabet order without changing its size: loading must fail before model state
is mutated. Try saving after backward but before the update: the writer must
reject pending gradients. Finally, list the state you would add for a scheduler,
dropout, mixed-precision scaling and two distributed workers.

\checkpoint
\begin{enumerate}
\item What does an atomic rename guarantee, and what does it not guarantee?
\item Why is a dataset cursor not sufficient for a shuffled or sampled data stream?
\item Why is equality of generated text a weaker conversion test than a full-logit comparison?
\item Which additional evidence would justify calling this a distributed, recoverable trainer?
\end{enumerate}
